\chapter{Separators and Divide-and-Conquer on Graphs}

Divide-and-conquer is the oldest trick in the algorithmist's book. Sort a list by splitting it in half; multiply integers by splitting their digits; build a balanced search tree by splitting at the median. Each time the move is the same: cut the input into pieces that are individually easier, recurse, and stitch the answers together. The cut is cheap (a midpoint), the pieces are balanced (each half the size), and the boundary between them is tiny (a single index). That is what makes the recursion $T(n)=2T(n/2)+\Oh(n)$ collapse to $\Oh(n\log n)$.

Graphs resist this. A graph is not a line, and there is usually no ``midpoint'' to split at. Worse, the boundary between two halves of a graph can be enormous: in the complete graph $K_n$, \emph{any} split into two halves of size $n/2$ leaves $n^2/4$ edges crossing between them. You cannot recurse usefully across a boundary that big. The whole game of divide-and-conquer on graphs comes down to a single question:

\begin{quote}
Can we split the vertices into two roughly equal halves by deleting only a \emph{small} set of vertices --- a set whose removal leaves no large piece behind?
\end{quote}

Such a small set is a \emph{separator}. When small balanced separators exist and we can find them, graph problems that look hopeless --- shortest paths through dense road networks, graph isomorphism, NP-hard problems on planar inputs --- crack open into recursions that actually terminate fast. The central result is the \emph{Planar Separator Theorem} of Lipton and Tarjan: every planar graph on $n$ vertices has a balanced separator of size $\Oh(\sqrt n)$. That single $\sqrt n$ is the seed of an entire algorithmic ecosystem, including sub-exponential $2^{\Oh(\sqrt n)}$ algorithms for problems that are otherwise believed to need $2^{\Om(n)}$.

This chapter develops separators from the ground up: what they are, why small ones are not guaranteed in general, how the planar theorem is proved (with its beautiful BFS-layering argument), how to wield it for divide-and-conquer, and how treewidth gives a second, parameter-driven source of separators. We close with the tradeoffs --- separator quality versus findability versus what the recursion actually buys you.

\chapterline

\section{What a separator is, and why balance matters}

\begin{definition}[Vertex separator]\label{ch27:def-sep}
Let $\Gph=(V,E)$ be an undirected graph and let $\alpha\in[\tfrac12,1)$ be a balance constant. A set $S\subseteq V$ is an \emph{$\alpha$-balanced (vertex) separator} if the vertices $V\setminus S$ can be partitioned into two sets $A,B$ with
\[
  A\cup B = V\setminus S,\qquad A\cap B=\emptyset,\qquad \text{no edge joins $A$ to $B$,}
\]
and $\abs{A}\le \alpha\,\abs{V}$, $\abs{B}\le \alpha\,\abs{V}$. The \emph{size} of the separator is $\abs S$.
\end{definition}

The two conditions pull in opposite directions. ``No edge joins $A$ to $B$'' means $S$ genuinely \emph{separates}: every path from $A$ to $B$ must pass through a deleted vertex. ``$\abs A,\abs B\le\alpha n$'' means the separation is \emph{balanced}: neither leftover piece is almost the whole graph. We want both, with $\abs S$ as small as possible.

\begin{intuitionbox}
Think of the graph drawn on the floor and imagine pouring a thin wall of removed vertices across it. The wall ($S$) cuts the floor into two regions ($A$ and $B$). A \emph{good} wall is short (few vertices) and splits the floor area fairly evenly. A long, snaking wall that fences off a single closet is balanced-but-useless if it is long; a one-brick wall that isolates a single dust bunny in the corner is short-but-useless because it is wildly unbalanced ($\abs A = 1$, $\abs B = n-2$). The separator theorem promises a wall that is both short and fair.
\end{intuitionbox}

Why insist on balance? Because divide-and-conquer's running time is governed by how fast the subproblems shrink. If we recurse on $A$ and $B$ and the larger of them can be as big as $n-1$, the recursion depth is $n$ and we have gained nothing. Balance ($\abs A,\abs B\le \tfrac23 n$, say) guarantees the larger child is a constant factor smaller, so the recursion has depth $\Oh(\log n)$. The separator size $\abs S$ controls how much ``boundary work'' each level costs. The two together --- small size and good balance --- are what make a separator theorem an \emph{algorithmic} theorem rather than a curiosity.

\begin{focusbox}
The three numbers that describe a separator family, and what each one buys:
\begin{itemize}
  \item \textbf{Size} $f(n)$: how many vertices you delete. Smaller is better; $\Oh(\sqrt n)$ for planar, $\Oh(1)$ (a single vertex) for trees, $\tw(\Gph)+1$ via treewidth.
  \item \textbf{Balance} $\alpha$: the larger leftover piece has at most $\alpha n$ vertices. We will mostly hit $\alpha=\tfrac23$; the exact constant rarely matters, only that $\alpha<1$.
  \item \textbf{Recursiveness}: ideally the class is closed under taking subgraphs, so the children are again separable by the same theorem. Planar graphs and bounded-treewidth graphs both have this; it is what lets the recursion repeat all the way down.
\end{itemize}
\end{focusbox}

\subsection{Trees: the warm-up that already wins}

Before anything hard, the easiest separator theorem in graph theory. Recall the \emph{centroid} of a tree (Chapter~20): a vertex whose removal leaves every component with at most $n/2$ vertices.

\begin{theorem}[Tree separator]\label{ch27:thm-tree}
Every tree on $n\ge 1$ vertices has a single vertex $c$ whose deletion splits it into components each of size at most $\floor{n/2}$. Thus a tree has a $\tfrac12$-balanced separator of size $1$.
\end{theorem}

\begin{proof}
Root the tree anywhere and let $\textit{sz}(v)$ be its subtree size. Start at the root and repeatedly walk to the child whose subtree is largest, stopping at the first vertex $c$ all of whose children $u$ satisfy $\textit{sz}(u)\le n/2$. At $c$, every ``downward'' component (a child subtree) has size $\le n/2$ by the stopping rule, and the single ``upward'' component has size $n-\textit{sz}(c)$. We claim $n-\textit{sz}(c)\le n/2$. We moved to $c$ from its parent because $\textit{sz}(c)>n/2$ (otherwise we would have stopped earlier, at the parent), so $n-\textit{sz}(c)<n/2$. Hence all components are $\le n/2$.
\end{proof}

This already powers \emph{centroid decomposition}: split at the centroid, recurse on each component, get a recursion tree of depth $\Oh(\log n)$. Trees are the friendliest case because a single vertex always suffices. The question this chapter answers is how far that idea reaches when the graph has cycles.

\subsection{The bad news: general graphs have no small separators}

You cannot hope for $\Oh(\sqrt n)$ separators in every graph. The obstruction is \emph{expansion} --- the very property we will celebrate in the chapters on expanders (Chapter~34). For now, the simplest witness:

\begin{proposition}[Cliques and expanders resist separation]\label{ch27:prop-clique}
In the complete graph $K_n$, any $\alpha$-balanced separator with $\alpha<1$ has size $\abs S\ge (1-\alpha)n=\Om(n)$: since every two vertices are adjacent, $A$ and $B$ cannot both be nonempty, so one side is empty and $\abs S \ge n - \max(\abs A,\abs B)\ge n-\alpha n$.
\end{proposition}

\begin{proof}
In $K_n$ every pair of vertices is joined by an edge. If $A$ and $B$ are both nonempty, pick $a\in A,b\in B$; the edge $ab$ joins $A$ to $B$, contradicting the no-edge condition. So one side, say $B$, is empty, forcing $A\cup S=V$ and $\abs S = n-\abs A\ge n-\alpha n=(1-\alpha)n$.
\end{proof}

A clique is dense, but density is not the real issue --- \emph{expansion} is. There exist sparse graphs (constant degree, only $\Oh(n)$ edges) called \emph{expanders} in which every balanced separator still has size $\Om(n)$. Removing a sublinear set of vertices from an expander always leaves one giant connected lump. So the planar separator theorem is not a free lunch you could have expected for all sparse graphs; it is a genuine structural payoff of planarity.

\begin{historybox}
The notion that ``nice'' graphs split well goes back to the 1970s VLSI and nested-dissection literature. George (1973) used $\sqrt n$-sized separators of the $\sqrt n\times\sqrt n$ grid to order Gaussian elimination so that fill-in stayed small --- \emph{nested dissection}. Lipton and Tarjan (1979) proved the general planar separator theorem, showing every planar graph behaves like a grid in this respect, and immediately spun out a string of applications: approximation, the nested-dissection bound for sparse linear systems, and lower bounds for Boolean circuits. Gilbert, Hutchinson, and Tarjan, and later Alon, Seymour, and Thomas, sharpened the constant. The expander obstruction was crystallized through the 1980s work on edge-expansion and eigenvalues (Cheeger, Alon--Milman), which Chapter~34 takes up.
\end{historybox}

\section{A grid by hand: the canonical $\sqrt n$ separator}

The single most clarifying example is the $k\times k$ grid graph $\Gph_{k,k}$, with $n=k^2$ vertices arranged in a square lattice and edges between orthogonal neighbours. Here you can \emph{see} the $\sqrt n$ separator with your eyes.

\begin{figure}[H]
\centering
\begin{tikzpicture}[scale=0.95,
  gv/.style={circle,draw,fill=blue!8,minimum size=5.5mm,inner sep=0pt,font=\scriptsize},
  sv/.style={circle,draw=trapred,fill=trapred!18,line width=0.7pt,minimum size=5.5mm,inner sep=0pt,font=\scriptsize},
  thick]
  \foreach \r in {0,...,4}{
    \foreach \c in {0,...,4}{
      \pgfmathtruncatemacro{\isS}{ifthenelse(\c==2,1,0)}
      \ifnum\isS=1
        \node[sv] (n\r\c) at (\c*1.25,\r*1.25) {};
      \else
        \node[gv] (n\r\c) at (\c*1.25,\r*1.25) {};
      \fi
    }
  }
  \foreach \r in {0,...,4}{
    \foreach \c in {0,...,3}{
      \pgfmathtruncatemacro{\cn}{\c+1}
      \draw[edgegray] (n\r\c)--(n\r\cn);
    }
  }
  \foreach \c in {0,...,4}{
    \foreach \r in {0,...,3}{
      \pgfmathtruncatemacro{\rn}{\r+1}
      \draw[edgegray] (n\r\c)--(n\rn\c);
    }
  }
  \node[font=\small] at (0.6,-1) {$A$ (left, $10$ vertices)};
  \node[font=\small,trapred] at (2.5,-1) {$S$ (middle column)};
  \node[font=\small] at (4.4,-1) {$B$ (right, $10$ vertices)};
\end{tikzpicture}
\caption{The $5\times 5$ grid, $n=25$. Deleting the middle column $S$ (red, $5=\sqrt n$ vertices) leaves $A$ (two left columns, $10$ vertices) and $B$ (two right columns, $10$ vertices) with no edge between them. A perfectly balanced separator of size $\sqrt n$.}
\label{ch27:fig-grid}
\end{figure}

\begin{labbox}
Verify Figure~\ref{ch27:fig-grid} by hand. The middle column has $k=5$ vertices. Every horizontal edge that would cross from a left column to a right column must pass through the middle column --- there is no edge skipping over it, because grid edges only join immediate neighbours. So removing the middle column severs all left--right connections. The leftover sides have $2k=10$ vertices each, both $\le \tfrac{2}{5} n$. Now ask: could you do it with fewer than $k$ deleted vertices? Intuitively no --- any wall separating left from right must block all $k$ rows, and the rows are vertex-disjoint paths, so by Menger's theorem (Chapter~12) you need at least $k$ vertices. The grid's separator number is exactly $\sqrt n$.
\end{labbox}

\begin{whybox}
Why is the grid the ``hardest easy'' case? Because it is the densest planar mesh: it has $\Th(n)$ edges and uniform local connectivity, so it looks locally like an expander. Yet globally it has a $\sqrt n$ separator. The planar separator theorem says this is no accident of the grid's symmetry --- \emph{every} planar graph, however irregular, can be cut as cleanly as a grid (up to constants). The grid is the extremal example that fixes the $\sqrt n$ scaling: you cannot beat $\sqrt n$ in general, because the grid already needs that many.
\end{whybox}

The grid also explains the recursion. A $\sqrt n$ separator splits $n$ vertices into two halves of $\le \tfrac23 n$. Recurse. The boundary work at the top level is $\Oh(\sqrt n)$, at the next level $\Oh(\sqrt{2n/3})$ on each of two pieces, and so on. The geometric decay of $\sqrt{\cdot}$ down a depth-$\Oh(\log n)$ tree is exactly what makes separator-based recurrences sum to small totals --- we make this precise in Section~\ref{ch27:sec-recur}.

\section{The Planar Separator Theorem}

\begin{theorem}[Lipton--Tarjan, 1979]\label{ch27:thm-lt}
Let $\Gph=(V,E)$ be a planar graph on $n$ vertices. Then there is a partition of $V$ into three sets $A$, $S$, $B$ such that
\begin{enumerate}
  \item no edge joins a vertex of $A$ to a vertex of $B$;
  \item $\abs A\le \tfrac23 n$ and $\abs B\le \tfrac23 n$;
  \item $\abs S \le 2\sqrt 2\,\sqrt n \approx 2.83\sqrt n$.
\end{enumerate}
Moreover $A$, $S$, $B$ can be computed in $\Oh(n)$ time.
\end{theorem}

The proof is constructive and is, at heart, a clever use of breadth-first search. We build up to it in three movements: (i) reduce to triangulated graphs, (ii) find a good \emph{level} of a BFS to use as the separator backbone, (iii) if no single level is small enough, cut the graph between two levels and patch with a spanning-tree path. The argument is one of the most elegant in algorithmic graph theory; we give it in enough detail to implement.

\begin{intuitionbox}
The idea in one breath: do a BFS from any root. This sorts the vertices into \emph{levels} $L_0,L_1,L_2,\dots$ by distance from the root. Each level is a ``moat'' --- removing all vertices of one level disconnects the inside (smaller-numbered levels) from the outside (larger-numbered levels), because every edge of a graph joins vertices in the same or adjacent levels. So \emph{any single BFS level is automatically a separator}; the only question is whether some level is both small ($\Oh(\sqrt n)$) and well-placed (balanced). Most of the proof is bookkeeping to guarantee such a level exists, using planarity (via Euler's formula) to bound how the levels can grow.
\end{intuitionbox}

\subsection{Movement 0: reductions}

We may assume $\Gph$ is connected (else separate the largest component using the rest as ballast; standard packing makes the pieces balanced). We may also assume each vertex has a nonnegative \emph{weight} summing to $1$, and ask only that each of $A,B$ carries weight $\le \tfrac23$; the unit-weight case is what we stated. Finally we may assume $\Gph$ is \emph{triangulated} (a \emph{maximal} planar graph, every face a triangle): adding edges only makes separation harder, and a triangulation can be computed in linear time. From here on $\Gph$ is a connected triangulated planar graph with $m\le 3n-6$ edges (Euler).

\subsection{Movement 1: BFS levels are separators}

Run BFS from an arbitrary root $r$. Let $L_i=\{v:\dist(r,v)=i\}$ be the $i$-th level and let $\ell$ be the last nonempty level. Write $L_{<i}=\bigcup_{j<i}L_j$ and $L_{>i}=\bigcup_{j>i}L_j$.

\begin{lemma}[Levels separate]\label{ch27:lem-levelsep}
For every $i$, the level $L_i$ is a separator: there is no edge between $L_{<i}$ and $L_{>i}$.
\end{lemma}

\begin{proof}
An edge $\{u,v\}$ of $\Gph$ joins vertices whose BFS distances differ by at most $1$ (BFS distances are a metric and adjacent vertices are within $1$). So if $u\in L_{<i}$, i.e.\ $\dist(r,u)\le i-1$, then $\dist(r,v)\le i$, hence $v\in L_{<i}\cup L_i$, never in $L_{>i}$.
\end{proof}

So if some level $L_i$ has $\abs{L_i}\le c\sqrt n$ \emph{and} both $\abs{L_{<i}}$ and $\abs{L_{>i}}$ are $\le \tfrac23 n$, we are done with $S=L_i$. Pick the \emph{middle level} $L_{i^*}$ at the median vertex: let $i^*$ be the smallest $i$ with $\abs{L_{\le i}}\ge n/2$. Then both sides have $\le n/2 \le \tfrac23 n$ vertices, so balance is automatic. The only thing that can go wrong is that this particular level is \emph{fat} --- it could have many more than $\sqrt n$ vertices.

The picture to keep in mind is a set of concentric moats around the root. Every edge stays inside one moat or hops to the next; none leaps across. Figure~\ref{ch27:fig-bfs} draws the levels of the $5\times5$ grid rooted at a corner, with the median level shaded.

\begin{figure}[H]
\centering
\begin{tikzpicture}[scale=0.92,
  l0/.style={circle,draw,fill=blue!18,minimum size=6mm,inner sep=0pt,font=\scriptsize},
  med/.style={circle,draw=trapred,fill=trapred!22,line width=0.7pt,minimum size=6mm,inner sep=0pt,font=\scriptsize},
  reg/.style={circle,draw,fill=blue!6,minimum size=6mm,inner sep=0pt,font=\scriptsize},
  thick]
  \foreach \r in {0,...,4}{
    \foreach \c in {0,...,4}{
      \pgfmathtruncatemacro{\lev}{\r+\c}
      \ifnum\lev=0
        \node[l0] (n\r\c) at (\c*1.3,\r*1.3) {\lev};
      \else
        \ifnum\lev=4
          \node[med] (n\r\c) at (\c*1.3,\r*1.3) {\lev};
        \else
          \node[reg] (n\r\c) at (\c*1.3,\r*1.3) {\lev};
        \fi
      \fi
    }
  }
  \foreach \r in {0,...,4}{
    \foreach \c in {0,...,3}{
      \pgfmathtruncatemacro{\cn}{\c+1}
      \draw[edgegray] (n\r\c)--(n\r\cn);
    }
  }
  \foreach \c in {0,...,4}{
    \foreach \r in {0,...,3}{
      \pgfmathtruncatemacro{\rn}{\r+1}
      \draw[edgegray] (n\r\c)--(n\rn\c);
    }
  }
  \node[font=\small] at (2.6,-1.0) {root at lower-left; label $=\dist(r,\cdot)=r+c$};
\end{tikzpicture}
\caption{BFS levels of the $5\times5$ grid rooted at the lower-left corner. Each vertex is labelled with its level $\dist(r,\cdot)$; on this grid the level is just the anti-diagonal index $r+c$, running from $0$ to $8$. Level sizes are $1,2,3,4,\mathbf 5,4,3,2,1$. The median level $L_4$ (red, the long anti-diagonal) holds $5=\sqrt n$ vertices and splits the $25$ vertices into $10$ below and $10$ above --- a thin median, so Movement~1 alone already separates the grid.}
\label{ch27:fig-bfs}
\end{figure}

\begin{labbox}
Read the level sizes straight off Figure~\ref{ch27:fig-bfs} and watch the proof run. The cumulative counts are $1,3,6,10,15,19,22,24,25$. The first cumulative count to reach $n/2=12.5$ is at level $4$ (cumulative $15$), so $i^*=4$: the median level is the long anti-diagonal $L_4=\{(4,0),(3,1),(2,2),(1,3),(0,4)\}$, exactly $\sqrt n$ vertices. Below it sit levels $0,1,2,3$ with $1+2+3+4=10$ vertices; above it sit levels $5,6,7,8$ with $4+3+2+1=10$. No edge of the grid joins a $\le 3$ vertex to a $\ge 5$ vertex --- check any edge and its two endpoints differ in level by exactly $1$. So $S=L_4$ is a balanced separator of size $\sqrt n$, and the grid never needs the harder Movements~2--3. The fat-median case only bites when the level-size profile spikes in the middle instead of forming this gentle tent; Exercise~\ref{ch27:ex-break} builds such a spike on purpose.
\end{labbox}

\subsection{Movement 2: a thin level must exist near the middle}

Here is where $\sqrt n$ enters. Suppose the median level $L_{i^*}$ is fat. We look for a \emph{thin} level above and a thin level below to sandwich it.

\begin{lemma}[A thin level nearby]\label{ch27:lem-thin}
There exist levels $a\le i^*\le b$ with $\abs{L_a}\le \sqrt n$ and $\abs{L_b}\le \sqrt n$, with $a$ within $\sqrt n$ levels above $i^*$ and $b$ within $\sqrt n$ levels below, such that cutting at $L_a$ and $L_b$ isolates the ``middle band'' $M=\bigcup_{a<j<b}L_j$ from the rest. In particular $\abs{L_a}+\abs{L_b}\le 2\sqrt n$ and $b-a\le 2\sqrt n$.
\end{lemma}

\begin{proof}
Walk upward from $i^*$ toward the root. If every one of the levels $L_{i^*}, L_{i^*-1},\dots, L_{i^*-\sqrt n}$ had more than $\sqrt n$ vertices, those $\sqrt n+1$ disjoint levels would together contain more than $\sqrt n\cdot\sqrt n=n$ vertices, impossible. So within $\sqrt n$ steps upward there is a level $L_a$ with $\abs{L_a}\le\sqrt n$. The same averaging argument applied downward from $i^*$ produces a thin level $L_b\le\sqrt n$ within $\sqrt n$ steps below. Hence $b-a\le 2\sqrt n$ and $\abs{L_a}+\abs{L_b}\le 2\sqrt n$.
\end{proof}

After Lemma~\ref{ch27:lem-thin} we have three regions: the \emph{top} $T=L_{<a}$, the \emph{middle band} $M=\bigcup_{a<j<b}L_j$, and the \emph{bottom} $B'=L_{>b}$, separated by the two thin levels $L_a$ and $L_b$ (total $\le 2\sqrt n$). If the middle band $M$ has $\abs M\le \tfrac23 n$, set $S=L_a\cup L_b$, $A=M$, $B=T\cup B'$. The top $T$ and bottom $B'$ are not adjacent (they are separated by $M$ and by the thin levels), so packing them together on one side is legal, and both sides are $\le\tfrac23 n$. Done, with $\abs S\le 2\sqrt n$.

\subsection{Movement 3: a fat middle band, cut by a tree path}

The remaining case is a middle band $M$ with more than $\tfrac23 n$ vertices but spanning only ``few'' levels (at most $b-a\le 2\sqrt n$ of them). Now planarity does the heavy lifting. Contract everything above $L_a$ to a single root and discard everything below $L_b$; the band $M$ together with $L_a,L_b$ is planar and triangulated, and its BFS depth from the contracted root is at most the number of band levels, $h\le 2\sqrt n$.

\begin{lemma}[Fundamental-cycle separator]\label{ch27:lem-cycle}
In a triangulated planar graph with a rooted spanning tree of height $h$, there is a non-tree edge whose induced \emph{fundamental cycle} (the tree path between its endpoints plus the edge) is a separator with both sides carrying $\le\tfrac23$ of the weight. The cycle has length $\le 2h+1$.
\end{lemma}

\begin{proof}[Idea]
Each non-tree edge $e$ closes a fundamental cycle $C_e$ with the tree (Figure~\ref{ch27:fig-cycle}). In a planar embedding, $C_e$ is a closed curve and, by the Jordan curve theorem, splits the plane --- hence the faces, hence the vertices --- into an inside and an outside. Walk over the faces: because the graph is triangulated, moving the dividing cycle across one triangular face changes the inside-weight by a bounded amount, so some non-tree edge yields a cycle with inside- and outside-weight both $\le\tfrac23 n$. Its vertex set is the separator, of size $\le 2h+1$, since the cycle consists of two root-to-endpoint tree paths (each $\le h$) plus the single non-tree edge.
\end{proof}

Apply Lemma~\ref{ch27:lem-cycle} inside the band. The effective height is $h\le 2\sqrt n$, so the cycle separator has $\le 2(2\sqrt n)+1=\Oh(\sqrt n)$ vertices. Add the two thin levels $L_a,L_b$ already counted; the grand total stays $\Oh(\sqrt n)$, and the original paper's careful accounting yields $\abs S\le 2\sqrt 2\sqrt n$. Every step --- BFS, level counting, choosing the thin levels, the fundamental-cycle search --- is linear time, proving the $\Oh(n)$ bound.

\begin{figure}[H]
\centering
\begin{tikzpicture}[scale=1.0,thick,
  tv/.style={circle,draw,fill=blue!8,minimum size=5mm,inner sep=0pt,font=\scriptsize},
  cv/.style={circle,draw=trapred,fill=trapred!18,line width=0.7pt,minimum size=5mm,inner sep=0pt,font=\scriptsize}]
  % root and spanning tree
  \node[tv] (r)  at (0,2.6) {$r$};
  \node[cv] (u)  at (-1.6,1.3) {$u$};
  \node[cv] (w)  at (1.6,1.3) {$w$};
  \node[tv] (p)  at (-2.4,0) {};
  \node[tv] (q)  at (2.4,0) {};
  \node[tv] (z)  at (0,0) {};
  % tree edges (solid)
  \draw[edgegray] (r)--(u); \draw[edgegray] (r)--(w);
  \draw[edgegray] (u)--(p); \draw[edgegray] (w)--(q);
  \draw[edgegray] (u)--(z); \draw[edgegray] (w)--(z);
  % non-tree edge e=uw (dashed), closing the fundamental cycle r-u-w-r
  \draw[trapred,dashed,line width=1pt] (u)--(w) node[midway,above,font=\scriptsize]{$e$ (non-tree)};
  % shade inside
  \begin{scope}[on background layer]
    \fill[trapred!7] (r.center)--(u.center)--(w.center)--cycle;
  \end{scope}
  \node[font=\scriptsize] at (0,1.75) {inside};
  \node[font=\scriptsize] at (-2.7,2.0) {outside};
\end{tikzpicture}
\caption{A fundamental-cycle separator. The non-tree edge $e=\{u,w\}$ (dashed red) together with the two tree paths from $u$ and $w$ up to their meeting point $r$ forms a cycle $C_e=r\,u\,w\,r$ (red). In the plane drawing $C_e$ bounds an inside and an outside (Jordan curve theorem); its vertices $\{r,u,w\}$ are the separator. Sweeping $e$ over triangular faces changes the inside-weight by a bounded amount, so some choice of $e$ balances the two sides to $\le\tfrac23 n$ each. The cycle length is at most two tree-path lengths plus one, i.e.\ $\le 2h+1$.}
\label{ch27:fig-cycle}
\end{figure}

\begin{whybox}
Where exactly was planarity used? Twice. First, in the edge bound $m\le 3n-6$ from Euler's formula, which forces BFS levels to grow ``slowly'' on average --- in a general graph a single BFS level could be all of $V\setminus\{r\}$ (a star), with no thin neighbour. Second, and more deeply, in Lemma~\ref{ch27:lem-cycle}: a cycle in a \emph{plane} drawing has a well-defined inside and outside (the Jordan curve theorem), which is what lets a single fundamental cycle act as a separator. Take away the plane embedding and a cycle separates nothing.
\end{whybox}

\begin{trapbox}
A tempting misreading: ``the smallest BFS level is a balanced separator.'' False. The smallest level might be $L_\ell$, a single far-flung vertex --- size $1$, but wildly unbalanced ($A=$ everything, $B=\emptyset$). Balance must be enforced by choosing the level near the \emph{median} of the vertex count, not the level of minimum size. The whole subtlety of Movements 2--3 is reconciling ``small'' with ``balanced'' when the median level happens to be fat. Confusing the two is the classic first-attempt bug.
\end{trapbox}

\begin{historybox}
Lipton and Tarjan's 1979 proof is the one above (with the constant $2\sqrt 2$). Djidjev (1982) improved the constant toward $\sqrt 6\approx 2.45$ and below. The Alon--Seymour--Thomas (1990) machinery and later cycle-separator refinements give strong guarantees. For the closely related notion of \emph{$r$-divisions} --- recursively dividing a planar graph into $\Oh(n/r)$ pieces of size $\Oh(r)$ each with $\Oh(\sqrt r)$ boundary --- Frederickson (1987) gave the construction underlying nearly every fast planar shortest-path algorithm. The whole subject generalizes: any graph excluding a fixed minor $H$ has an $\Oh(\sqrt n)$ separator (Alon--Seymour--Thomas), so the theorem is really about minor-closed families, of which planar graphs are the flagship.
\end{historybox}

\begin{algorithm}[H]
\caption{\textsc{PlanarSeparator}$(\Gph)$ --- Lipton--Tarjan outline}
\label{ch27:alg-sep}
\DontPrintSemicolon
\KwIn{connected planar graph $\Gph=(V,E)$, $n=\abs V$}
\KwOut{separator $S$, parts $A$, $B$ with $\abs A,\abs B\le\tfrac23 n$, $\abs S=\Oh(\sqrt n)$}
Triangulate $\Gph$ (linear time); fix a planar embedding\;
Run BFS from an arbitrary root $r$; record levels $L_0,L_1,\dots,L_\ell$\;
Let $i^*$ be the least $i$ with $\abs{L_{\le i}}\ge n/2$ \tcp*{median level}
\If{$\abs{L_{i^*}}\le 2\sqrt n$}{
  \Return $S\gets L_{i^*}$, $A\gets L_{<i^*}$, $B\gets L_{>i^*}$ \tcp*{thin median: done}
}
Find $a\le i^*$ (nearest) with $\abs{L_a}\le\sqrt n$\;
Find $b\ge i^*$ (nearest) with $\abs{L_b}\le\sqrt n$ \tcp*{thin sandwich, $b-a=\Oh(\sqrt n)$}
$T\gets L_{<a}$;\quad $M\gets L_{(a,b)}$;\quad $B'\gets L_{>b}$\;
\If{$\abs M\le \tfrac23 n$}{
  \Return $S\gets L_a\cup L_b$, $A\gets M$, $B\gets T\cup B'$
}
Within the band $M\cup L_a\cup L_b$ build a BFS tree and run \textsc{FundamentalCycleSep} (Lemma~\ref{ch27:lem-cycle})\;
\Return cycle $C$ unioned with $L_a\cup L_b$ as $S$; balanced face-partition gives $A,B$\;
\end{algorithm}

\section{Separator-driven recurrences: what the $\sqrt n$ buys}
\label{ch27:sec-recur}

Now the payoff. A separator theorem turns a problem on $\Gph$ into the same problem on two graphs each $\le\tfrac23 n$ vertices, plus boundary work proportional to $\abs S$. The exact dividend depends on how the boundary interacts with the recursion. Three recurrence shapes recur (pun intended).

\paragraph{Linear divide-and-conquer.} If recombining costs $\Oh(n)$ and the separator is $\Oh(\sqrt n)$,
\[
  T(n) \le T(n_1)+T(n_2)+\Oh(n),\qquad n_1+n_2\le n,\ n_1,n_2\le\tfrac23 n.
\]
This solves to $T(n)=\Oh(n\log n)$, exactly like mergesort: the work per level is $\Oh(n)$ (summed across all pieces at that depth) and there are $\Oh(\log n)$ levels. The $\sqrt n$ does not even show up --- it is dominated by the $\Oh(n)$ merge.

\paragraph{Boundary-dominated work.} When the cost at a node is proportional to the \emph{separator size} rather than the subproblem size --- common when the separator vertices become ``interface'' variables --- the recurrence is
\[
  T(n) \le T(n_1)+T(n_2)+\Oh(\sqrt n).
\]
The key inequality is concavity of the square root:
\[
  \sqrt{n_1}+\sqrt{n_2}\le \sqrt2\,\sqrt n \quad\text{when } n_1+n_2=n,
\]
so each level multiplies the total boundary by at most $\sqrt2$, while the pieces shrink by a constant factor. Over $\Oh(\log n)$ levels this geometric tug-of-war resolves to $\Oh(\sqrt n\log n)$ for the unbalanced split and $\Oh(\sqrt n)$ when the recursion is stopped early (as in $r$-divisions).

\begin{whybox}
The recurrence $T(n)=T(2n/3)+T(n/3)+\Oh(\sqrt n)$ is worth solving once by hand. Guess $T(n)\le c\sqrt n$. Then $T(n)\le c\sqrt{2n/3}+c\sqrt{n/3}+\Oh(\sqrt n)=c\sqrt n(\sqrt{2/3}+\sqrt{1/3})+\Oh(\sqrt n)$. Now $\sqrt{2/3}+\sqrt{1/3}=0.816+0.577=1.394>1$, so the homogeneous part \emph{grows}; the guess $c\sqrt n$ fails and the true solution carries a $\log$ factor: $T(n)=\Oh(\sqrt n\log n)$. The lesson: $\sqrt n$ separators do not magically give $\Oh(\sqrt n)$ total work --- you usually pay one or two log factors, and you should always check the constant $\sqrt{\alpha}+\sqrt{1-\alpha}$ against $1$.
\end{whybox}

\paragraph{Sub-exponential exact algorithms.} The recurrence that makes separators famous appears in \emph{exponential-time} algorithms. Suppose solving a problem (graph colouring, independent set, Hamiltonicity) on $\Gph$ means trying all assignments to the separator $S$ and recursing. With $\abs S=\Oh(\sqrt n)$ and $q$ choices per separator vertex, the recurrence is
\[
  T(n)\le q^{\,c\sqrt n}\cdot\big(T(n_1)+T(n_2)\big).
\]
We solve it in Section~\ref{ch27:sec-subexp} and find $T(n)=2^{\Oh(\sqrt n)}$ --- sub-exponential, a dramatic speedup over the $2^{\Th(n)}$ brute force. This is the structural reason planar NP-hard problems are so much easier than general ones.

\begin{complexitybox}
Quick reference for separator recurrences ($\alpha$-balanced, separator size $f(n)$):
\begin{center}
\begin{tabular}{lll}
\toprule
Per-node cost & Recurrence & Solution \\
\midrule
$\Oh(n)$ merge & $T(n_1)+T(n_2)+\Oh(n)$ & $\Oh(n\log n)$ \\
$\Oh(\sqrt n)$ boundary & $T(\tfrac23 n)+T(\tfrac13 n)+\Oh(\sqrt n)$ & $\Oh(\sqrt n\log n)$ \\
$2^{\Oh(\sqrt n)}$ enumerate & $2^{c\sqrt n}\,[T(n_1)+T(n_2)]$ & $2^{\Oh(\sqrt n)}$ \\
treewidth $w$ DP & per bag $q^{w}$, $\Oh(n)$ bags & $\Oh(q^{w}\,\poly(w)\,n)$ \\
\bottomrule
\end{tabular}
\end{center}
The rule of thumb: a size-$f(n)$ separator buys you a $2^{\Oh(f(n))}$ exact algorithm and an $\Oh(f(n))$-factor in divide-and-conquer.
\end{complexitybox}

\section{Application 1: planar shortest paths and nested dissection}

The earliest and still most important use of planar separators is \emph{nested dissection} --- ordering the vertices so that recursive elimination touches as little of the graph as possible.

\subsection{Nested dissection for sparse linear systems}

To solve $\Lap x=b$ for a planar graph Laplacian $\Lap$ (Chapter~31), Gaussian elimination's cost is dominated by \emph{fill-in}: zero entries that become nonzero as you eliminate. Nested dissection eliminates the two halves $A,B$ \emph{before} the separator $S$. Because $A$ and $B$ share no edges, eliminating an $A$-vertex never creates fill between $A$ and $B$ --- the separator $S$ shields them. Recursing, the fill-in obeys a separator recurrence and one proves:

\begin{theorem}[George 1973; Lipton--Rose--Tarjan 1979]\label{ch27:thm-nd}
A planar Laplacian system can be solved by nested dissection with $\Oh(n\log n)$ fill-in and $\Oh(n^{3/2})$ arithmetic operations, versus $\Oh(n^3)$ for naive elimination on a dense matrix.
\end{theorem}

\begin{connectionbox}
This is the bridge to the LINEAR ALGEBRA companion and to Chapter~31 (Laplacians) and Chapter~32 (spectra). Nested dissection is why ``solve a planar Laplacian'' is genuinely cheaper than ``solve a general Laplacian'': the geometry of the graph controls the algebra of the matrix. The same separator structure underlies the fast \emph{Laplacian solvers} of Spielman--Teng and successors, which for planar inputs can be specialized using exactly these $\sqrt n$ separators. When you read about electrical flows (Chapter~35) on planar networks, the $\Oh(n^{3/2})$ here is the elimination-based ancestor of those near-linear solvers.
\end{connectionbox}

\subsection{Planar single-source shortest paths}

For shortest paths, separators give the celebrated $\Oh(n)$-time planar SSSP of Henzinger, Klein, Rao, and Subramanian, and more practically the \emph{$r$-division} approach to all-pairs queries and distance oracles. The skeleton:

\begin{enumerate}
  \item Compute an $r$-division: $\Oh(n/r)$ pieces, each $\le r$ vertices, each with a \emph{boundary} of $\Oh(\sqrt r)$ separator vertices shared with neighbours.
  \item Inside each piece, precompute shortest paths between its boundary vertices --- a \emph{dense distance graph} on $\Oh(\sqrt r)$ nodes per piece.
  \item Glue the dense distance graphs across the global boundary (total $\Oh(n/\sqrt r)$ boundary vertices) and run a shortest-path computation on this much smaller \emph{condensed} graph.
\end{enumerate}

\begin{intuitionbox}
You are planning a road trip across a continent divided into provinces. Within each province you precompute the driving time between its border crossings only --- you do not care about the interior towns once you are passing through. Then you plan the cross-country route purely on the network of border crossings. Because each province has only $\Oh(\sqrt r)$ crossings (its separator), the inter-province network is tiny compared to the full road map, and the global plan is cheap. The separator is precisely the set of border crossings; planarity guarantees each province has few of them.
\end{intuitionbox}

\begin{applicationbox}
This is not theoretical folklore. Real route planners on road networks (which are nearly planar) use exactly this hierarchy: \emph{customizable route planning} and \emph{multi-level Dijkstra} build separator-based overlays so that a continental shortest-path query touches only the boundary networks, answering in microseconds. The separator is the algorithmic content; planarity (approximately, via low highway dimension) is why it is small.
\end{applicationbox}

\subsection{A tiny worked trace}

Take the $4\times 4$ grid ($n=16$), source $s$ at the top-left corner. Use column $3$ as the separator $S$ (four vertices, splitting columns $1,2$ from column $4$).

\begin{labbox}
Do it on paper. Label vertices $(r,c)$, $r,c\in\{1,2,3,4\}$, unit edge weights, source $(1,1)$.
\begin{enumerate}
  \item \textbf{Separator} $S=\{(1,3),(2,3),(3,3),(4,3)\}$ (column 3). Left piece $A=$ columns $1,2$ (8 vertices), right piece $B=$ column $4$ (4 vertices).
  \item \textbf{Inside $A$}: run Dijkstra from $s$. Distances to the four separator vertices of column 3 equal their grid Manhattan distance from the corner: $\dist((1,3))=2$, $\dist((2,3))=3$, $\dist((3,3))=4$, $\dist((4,3))=5$.
  \item \textbf{Across $S$}: the only way into $B$ (column 4) is through column 3. So $\dist((r,4)) = \min_{r'} \big(\dist((r',3)) + \abs{r-r'} + 1\big)$, taking the cheapest border crossing.
  \item For $(1,4)$: best is via $(1,3)$: $2+0+1=3$. For $(4,4)$: via $(4,3)$: $5+0+1=6$, or via $(1,3)$: $2+3+1=6$ --- a tie, both Manhattan-correct.
\end{enumerate}
You never needed interior-to-interior distances across the boundary: every $A$-to-$B$ shortest path is captured by the four border values. That is the separator paying off in miniature.
\end{labbox}

\section{Application 2: sub-exponential algorithms on planar graphs}
\label{ch27:sec-subexp}

Now the main theoretical payoff. Many NP-hard problems --- \emph{maximum independent set}, \emph{vertex cover}, \emph{$3$-colouring}, \emph{Hamiltonian cycle} --- are still NP-hard restricted to planar graphs. But on planar inputs they admit $2^{\Oh(\sqrt n)}$ algorithms, versus $2^{\Th(n)}$ in general. The separator is the entire reason.

\begin{theorem}[Sub-exponential planar algorithms]\label{ch27:thm-subexp}
Maximum independent set (and many similar problems) on an $n$-vertex planar graph can be solved exactly in time $2^{\Oh(\sqrt n)}\cdot\poly(n)$.
\end{theorem}

\begin{proof}
Find a separator $S$, $\abs S\le c\sqrt n$, splitting $V$ into $A,B$ with $\abs A,\abs B\le\tfrac23 n$. The optimal independent set $I^*$ interacts with $S$ in some way: $I^*\cap S$ is one of at most $2^{\abs S}=2^{c\sqrt n}$ subsets of $S$. \emph{Enumerate} all independent subsets $X\subseteq S$. For each such $X$:
\begin{itemize}
  \item Forbid the neighbours of $X$ in $A$ and in $B$ (they cannot be in $I^*$).
  \item Because no edge crosses $A$--$B$, the best extension into $A$ and the best extension into $B$ are \emph{independent subproblems}: recurse on $\Gph[A]$ and $\Gph[B]$ with the forbidden vertices removed.
  \item Total for this $X$: $\abs X + \textsc{MIS}(A') + \textsc{MIS}(B')$.
\end{itemize}
Take the best over all $X$. The recurrence is
\[
  T(n)\le 2^{c\sqrt n}\,\big(T(\tfrac23 n)+T(\tfrac13 n)\big)+\poly(n).
\]
Let $T(n)=2^{g(n)}$. The dominant term gives $g(n)\le c\sqrt n + g(\tfrac23 n)$ (the larger child dominates). Unrolling, $g(n)\le c\sqrt n\sum_{k\ge0}(\sqrt{2/3})^{k}=c\sqrt n\cdot\frac{1}{1-\sqrt{2/3}}=\Oh(\sqrt n)$, because $\sqrt{2/3}<1$ makes the geometric series converge. Hence $T(n)=2^{\Oh(\sqrt n)}$.
\end{proof}

\begin{whybox}
The geometric series is the crux, and it is where planarity earns its keep. The separator size shrinks by $\sqrt{2/3}\approx 0.816$ each level (since $n\mapsto\tfrac23 n$ and we take square roots), and a geometric series with ratio $<1$ sums to a constant times its first term $c\sqrt n$. If the separator were size $\Th(n)$ (as in a general or expander graph), each level would contribute $\Th(n)$ to the exponent and the sum would be $\Th(n\log n)$ --- worse than brute force. The sub-exponential bound is \emph{exactly} the statement that the separator-size series converges, which is \emph{exactly} the statement that separators are sublinear. No sublinear separator, no sub-exponential algorithm.
\end{whybox}

\begin{algorithm}[H]
\caption{\textsc{PlanarMIS}$(\Gph)$ --- sub-exponential max independent set}
\label{ch27:alg-mis}
\DontPrintSemicolon
\KwIn{planar graph $\Gph=(V,E)$, $n=\abs V$}
\KwOut{size of a maximum independent set}
\If{$n\le n_0$}{\Return brute force}
$(A,S,B)\gets$\textsc{PlanarSeparator}$(\Gph)$\;
$\textit{best}\gets 0$\;
\For{each independent subset $X\subseteq S$}{
  $A'\gets A\setminus N(X)$;\quad $B'\gets B\setminus N(X)$ \tcp*{delete forbidden neighbours}
  $v\gets \abs X + \textsc{PlanarMIS}(\Gph[A']) + \textsc{PlanarMIS}(\Gph[B'])$\;
  $\textit{best}\gets\max(\textit{best},v)$\;
}
\Return \textit{best}\;
\end{algorithm}

\begin{applicationbox}
The same template --- ``enumerate the separator, recurse on the independent halves'' --- gives $2^{\Oh(\sqrt n)}$ algorithms for planar Hamiltonicity, planar TSP, planar Steiner tree, planar dominating set, and more. It also drives \emph{Baker's technique} for PTASs on planar graphs (peel off BFS layers $\bmod\ k$ to get bounded-treewidth slabs, solve each by DP, glue) --- a separator-flavoured method we revisit in the approximation chapter (Chapter~38). The slogan ``planar problems are sub-exponential'' is one of the cleanest manifestations of structure beating brute force.
\end{applicationbox}

\section{Application 3: graph isomorphism on planar graphs}

Separators were among Lipton and Tarjan's original applications: a fast planar graph isomorphism test. General graph isomorphism has no known polynomial algorithm (Babai's 2015 quasi-polynomial bound is the frontier), but \emph{planar} isomorphism is solvable in linear time --- and the early algorithms went through separators with a divide-and-conquer that matches the two graphs separator-first.

\begin{intuitionbox}
To decide whether two planar maps are the same, find a small separator in each. The separator splits each map into a few balanced pieces. Try to match a separator of $\Gph_1$ to a separator of $\Gph_2$ (only $\Oh(\sqrt n)$ vertices, so few alignments survive when you also exploit the embedding), then recursively match the pieces. Because the pieces are a constant factor smaller and the boundary is small, the matching recursion terminates fast. Modern planar isomorphism uses canonical embeddings and tree decompositions instead, but the separator viewpoint is the historical and conceptual root.
\end{intuitionbox}

\begin{historybox}
Hopcroft and Tarjan (1972) gave the first linear-time planar isomorphism test using a recursive decomposition into triconnected components; Lipton--Tarjan separators offered an alternative divide-and-conquer. For the harder \emph{bounded-genus} and \emph{bounded-treewidth} classes, isomorphism is also polynomial, again via separator/decomposition methods (Bodlaender). The deep modern statement is that isomorphism is fixed-parameter tractable in treewidth and in genus --- all descendants of ``small separators make matching tractable.''
\end{historybox}

\section{Separators from treewidth}

The second great source of separators is \emph{treewidth} (Chapter~24). Where planarity gives $\Oh(\sqrt n)$ separators via geometry, a tree decomposition gives separators of size $\tw(\Gph)+1$ via structure --- and crucially, these are \emph{recursive by construction}.

\begin{definition}[Recap: tree decomposition]\label{ch27:def-td}
A \emph{tree decomposition} of $\Gph=(V,E)$ is a tree $\mathcal T$ whose nodes carry \emph{bags} $B_t\subseteq V$ such that (i) every vertex appears in some bag, (ii) every edge $\{u,v\}$ has both ends in a common bag, and (iii) for each vertex $v$, the bags containing $v$ form a connected subtree of $\mathcal T$. The \emph{width} is $\max_t\abs{B_t}-1$, and $\tw(\Gph)$ is the minimum width over all decompositions.
\end{definition}

The connectivity condition (iii) is exactly what makes every bag a separator.

\begin{figure}[H]
\centering
\begin{tikzpicture}[scale=0.95,thick,
  gv/.style={circle,draw,fill=blue!8,minimum size=6mm,inner sep=0pt,font=\scriptsize},
  bag/.style={rounded corners,draw=structurepurple,fill=structurepurple!8,inner sep=3pt,font=\scriptsize}]
  % small graph: a 6-cycle with a chord -> treewidth 2
  \begin{scope}[xshift=0cm]
    \node[gv] (a) at (0,2)   {a};
    \node[gv] (b) at (1.4,2.4) {b};
    \node[gv] (c) at (2.2,1.2) {c};
    \node[gv] (d) at (1.4,0)  {d};
    \node[gv] (e) at (0,0.4)  {e};
    \draw[edgegray] (a)--(b)--(c)--(d)--(e)--(a);
    \draw[edgegray] (b)--(e); \draw[edgegray] (c)--(e);
    \node[font=\small] at (1.1,-0.9) {graph $\Gph$};
  \end{scope}
  % tree decomposition, bags of size 3 => width 2
  \begin{scope}[xshift=5.4cm,yshift=0.2cm]
    \node[bag] (B1) at (0,2.2)  {$\{a,b,e\}$};
    \node[bag] (B2) at (0,1.0)  {$\{b,c,e\}$};
    \node[bag] (B3) at (0,-0.2) {$\{c,d,e\}$};
    \draw (B1)--(B2)--(B3);
    \node[font=\small] at (0,-1.1) {decomposition $\mathcal T$ (width 2)};
  \end{scope}
  \draw[->,structurepurple,dashed] (3.0,1.0) -- (4.0,1.0) node[midway,above,font=\scriptsize]{};
\end{tikzpicture}
\caption{A graph of treewidth $2$ and a width-$2$ tree decomposition. The shared bag $\{b,e\}$ between the top two bags is a separator: deleting $\{b,e\}$ disconnects $\{a\}$ from $\{c,d\}$. Every bag separates the graph into the parts hanging off the corresponding tree node.}
\label{ch27:fig-td}
\end{figure}

\begin{theorem}[Bags are separators]\label{ch27:thm-bagsep}
Let $\mathcal T$ be a tree decomposition of $\Gph$ of width $w$. For any tree edge $\{t,t'\}$ of $\mathcal T$, removing the bag intersection $B_t\cap B_{t'}$ separates the vertices ``on the $t$ side'' from those ``on the $t'$ side.'' In particular, every bag $B_t$ is a separator of size $\le w+1$ whose removal disconnects the graph into the parts corresponding to the components of $\mathcal T-t$.
\end{theorem}

\begin{proof}
Deleting tree edge $\{t,t'\}$ splits $\mathcal T$ into two subtrees, $\mathcal T_t$ and $\mathcal T_{t'}$. Let $A$ = vertices appearing only in bags of $\mathcal T_t$, $B$ = vertices appearing only in bags of $\mathcal T_{t'}$, and $S=B_t\cap B_{t'}$. Suppose for contradiction an edge $\{u,v\}\in E$ has $u\in A,v\in B$. By property (ii) some bag $B_x$ contains both $u$ and $v$. Say $x\in\mathcal T_t$. Then $v\in B_x$ with $x\in\mathcal T_t$, so $v$ appears in a bag of $\mathcal T_t$. But $v$ also appears in some bag of $\mathcal T_{t'}$ (as $v\in B$); by property (iii) the bags containing $v$ form a connected subtree, which therefore spans both sides of edge $\{t,t'\}$, forcing $v\in B_t\cap B_{t'}=S$ --- contradicting $v\in B$ (which excludes $S$). Hence no $A$--$B$ edge exists, and $S$ separates $A$ from $B$ with $\abs S\le\abs{B_t}\le w+1$.
\end{proof}

\begin{whybox}
Property (iii) --- the running-intersection / connected-subtree condition --- is the load-bearing wall. It says each vertex ``lives'' in a connected region of the tree, so to pass from a vertex living strictly on one side to one living strictly on the other, you must cross through a vertex living \emph{on the boundary bag}. That boundary bag is the separator. This is the same reason a single vertex separates a tree (Theorem~\ref{ch27:thm-tree}): a tree is exactly a width-$1$ decomposition of itself, and the centroid is the balanced bag.
\end{whybox}

\subsection{Balanced separators from bounded treewidth}

Theorem~\ref{ch27:thm-bagsep} gives \emph{some} separator; to get a \emph{balanced} one, find the centroid bag of the decomposition tree.

\begin{theorem}[Balanced separator from treewidth]\label{ch27:thm-tw-balanced}
A graph of treewidth $w$ has a $\tfrac12$-balanced separator $S$ with $\abs S\le w+1$: every component of $\Gph-S$ has $\le\tfrac12 n$ vertices.
\end{theorem}

\begin{proof}
Take a tree decomposition $\mathcal T$ of width $w$. Weight each tree node $t$ by the number of vertices appearing \emph{first} in $B_t$ when $\mathcal T$ is rooted (so the total weight is $n$ and each vertex is counted once). Find the centroid bag $t^*$ of $\mathcal T$ under these weights, as in Theorem~\ref{ch27:thm-tree}: removing $t^*$ from $\mathcal T$ leaves subtrees each carrying $\le\tfrac12 n$ weight. Set $S=B_{t^*}$. By Theorem~\ref{ch27:thm-bagsep}, $S$ separates the graph into the parts corresponding to those subtrees, each of size $\le\tfrac12 n$, and $\abs S\le w+1$.
\end{proof}

\begin{connectionbox}
This is the deep duality between the two halves of the divide-and-conquer story. \emph{Treewidth is small if and only if the graph has small balanced separators recursively.} One direction is Theorem~\ref{ch27:thm-tw-balanced} (small treewidth $\Rightarrow$ small balanced separators). The converse also holds up to constants: if every subgraph has a balanced separator of size $\le k$, then $\tw(\Gph)=\Oh(k\log n)$ (and $\Oh(k)$ for nice families). So ``has small separators'' and ``has small treewidth'' are essentially the same structural property viewed from two angles --- separators are the \emph{local} witness, a tree decomposition is the \emph{global} certificate assembling them. Planar graphs have $\tw=\Oh(\sqrt n)$ precisely because of the $\Oh(\sqrt n)$ separator theorem, and this is what makes the planar DP and FPT machinery of Chapters~24--26 run.
\end{connectionbox}

\begin{example}[Planar treewidth bound]\label{ch27:ex-planar-tw}
Applying the $\Oh(\sqrt n)$ planar separator recursively bounds planar treewidth: $\tw(\Gph)=\Oh(\sqrt n)$ for planar $\Gph$. The recursion $f(n)\le f(\tfrac23 n)+\Oh(\sqrt n)$ for the width of the assembled decomposition solves to $\Oh(\sqrt n)$. Combined with bounded-treewidth DP (cost $2^{\Oh(\tw)}\poly n$), this is a \emph{second} route to the $2^{\Oh(\sqrt n)}$ planar algorithms of Section~\ref{ch27:sec-subexp}: build a width-$\Oh(\sqrt n)$ decomposition, then run the standard treewidth DP. Same bound, structurally cleaner --- you separate once into a decomposition, then never separate again.
\end{example}

\subsection{Treewidth DP: bounded width gives exact poly-time}

When $w$ is a \emph{constant} (or small parameter), the separator is constant-sized and the recursion becomes a polynomial-time dynamic program over the bags --- the FPT engine of Chapter~25.

\begin{algorithm}[H]
\caption{\textsc{TreewidthDP} skeleton --- weighted independent set, width $w$}
\label{ch27:alg-twdp}
\DontPrintSemicolon
\KwIn{nice tree decomposition $(\mathcal T,\{B_t\})$ of width $w$, rooted}
\KwOut{max weight of an independent set}
\For{each node $t$ of $\mathcal T$ in post-order}{
  \For{each independent subset $X\subseteq B_t$}{
    \tcp{$\textit{dp}[t][X]$ = best weight in the subtree below $t$ with $B_t\cap I = X$}
    combine children's $\textit{dp}$ tables consistent with $X$ on the shared (separator) vertices\;
  }
}
\Return $\max_X \textit{dp}[\text{root}][X]$\;
\end{algorithm}

\begin{complexitybox}
The DP keeps, per bag, a table indexed by subsets of the bag --- $\Oh(2^{w})$ states for independent set, $2^{\Oh(w)}$ for connectivity problems like Hamiltonicity (via the Cut\&Count / rank-based methods of Chapter~25). With $\Oh(n)$ bags, total time is $\Oh(2^{w}\,w\,n)$ for independent set --- \emph{linear in $n$} for fixed $w$. The separator size $w+1$ is exactly the table's exponent: every join over a tree edge synchronizes on the shared separator $B_t\cap B_{t'}$, so the separator is what you enumerate, just as in the planar MIS algorithm. Treewidth DP \emph{is} separator-based divide-and-conquer, with the decomposition precomputing all the separators at once.
\end{complexitybox}

\begin{labbox}
Run the DP by hand on the graph of Figure~\ref{ch27:fig-td}, vertices $\{a,b,c,d,e\}$ with edges $ab,bc,cd,de,ea,be,ce$, all weights $1$, maximizing the size of an independent set. The non-edges are exactly $ac,ad,bd$, so the true answer is a $2$-set such as $\{a,c\}$, $\{a,d\}$, or $\{b,d\}$ --- no $3$-set is independent (check: any third vertex hits an edge). We will watch the DP rediscover $\mathbf 2$ touching only $\le w+1=3$ vertices at a time. Use the path decomposition $B_3=\{c,d,e\}$, $B_2=\{b,c,e\}$, $B_1=\{a,b,e\}$, rooted at $B_1$ so post-order is $B_3,B_2,B_1$. The shared separators are $B_2\cap B_3=\{c,e\}$ and $B_1\cap B_2=\{b,e\}$. Vertex $d$ is forgotten when we leave $B_3$; $c$ is forgotten when we leave $B_2$; $a$ is introduced only at the root. Each $\textit{dp}[B_t][X]$ stores the best independent set in the subtree at or below $t$ whose intersection with $B_t$ equals $X$.

\textbf{Leaf $B_3=\{c,d,e\}$.} Edges inside: $cd,de,ce$ --- a triangle, so only singletons and $\emptyset$ are independent. $\textit{dp}[B_3]$: $\textit{dp}[\emptyset]=0$, $\textit{dp}[\{c\}]=1$, $\textit{dp}[\{d\}]=1$, $\textit{dp}[\{e\}]=1$. Forgetting $d$, we project onto the separator $\{c,e\}$, keeping the best value per pattern: $\{c,e\}{=}\emptyset\to1$ (realized by $\{d\}$), $\{c\}\to1$, $\{e\}\to1$.

\textbf{Internal $B_2=\{b,c,e\}$.} Edges inside: $bc,ce,be$ --- again a triangle, only singletons. The forget-$d$ projection above tells us how much weight the subtree below contributes for each choice on the shared $\{c,e\}$: pattern $\emptyset\to1$ (the set $\{d\}$), $\{c\}\to1$ (the set $\{c\}$), $\{e\}\to1$ (the set $\{e\}$). For each independent $X\subseteq B_2$ we glue to the child pattern whose $\{c,e\}$ choice \emph{agrees} with $X$ (so we never combine states that disagree on a separator vertex), and we add only the weight of the vertices of $X$ \emph{not already counted in the child} (the shared separator vertices are counted once): $X=\emptyset\to 0+(\text{pattern }\emptyset)=0+1=1$; $X=\{b\}\to w(b)+(\text{pattern }\emptyset)=1+1=2$, since $b$ is new (it is not in the child's $\{c,e\}$); $X=\{c\}\to(\text{pattern }\{c\})=1$, the shared $c$ already counted in the child; $X=\{e\}\to(\text{pattern }\{e\})=1$, the shared $e$ already counted. Now forget $c$, projecting onto $\{b,e\}$ by taking, per pattern, the best value over $X$ that excludes $c$: $\emptyset\to1$ (from $X=\emptyset$, beating $X=\{c\}{=}1$), $\{b\}\to2$ (from $X=\{b\}$), $\{e\}\to1$ (from $X=\{e\}$).

\textbf{Root $B_1=\{a,b,e\}$.} Edges inside: $ab,be,ea$ --- triangle, only singletons. Vertex $a$ is introduced here; it is non-adjacent to the already-forgotten $c,d$, so adding $a$ never conflicts below. Combine each independent $X\subseteq B_1$ with the child's $\{b,e\}$-pattern that agrees with $X$, again adding only the weight of new (non-shared) vertices: $X=\emptyset\to 0+(\text{pattern }\emptyset)=0+1=1$; $X=\{a\}\to w(a)+(\text{pattern }\emptyset)=1+1=2$, with $a$ new and the child realizing pattern $\emptyset$ by $\{d\}$, so the set is $\{a,d\}$; $X=\{b\}\to(\text{pattern }\{b\})=2$, the shared $b$ already counted in the child pattern $\{b\}$ (realized by $\{b,d\}$), and $a\notin X$ so no $ab$ conflict arises; $X=\{e\}\to(\text{pattern }\{e\})=1$, shared $e$ counted once. The maximum over the root is $\mathbf 2$, realized by $\{a,d\}$ (from $X=\{a\}$, inheriting $d$ from $B_3$) or by $\{b,d\}$ (from $X=\{b\}$). The mechanism to take away: each join is a \emph{constrained} merge on the shared separator ($\{c,e\}$, then $\{b,e\}$) --- only states agreeing on every separator vertex may glue --- which is exactly why the separator size is the table's exponent and why bounded width gives linear-time exact answers.
\end{labbox}

\section{Tradeoffs: size, balance, findability, and what you actually gain}

Separator theorems are a menu, not a single dish. The choices interact, and picking the wrong point on the curve wastes the structure. The main axes:

\begin{itemize}
  \item \textbf{Size vs.\ balance.} You can often trade a smaller separator for worse balance or vice versa. A $\tfrac12$-balanced separator may be larger than a $\tfrac23$-balanced one. For divide-and-conquer the constant $\alpha$ barely matters (any $\alpha<1$ gives $\Oh(\log n)$ depth), so prefer the smallest separator with \emph{some} constant balance.
  \item \textbf{Vertex vs.\ edge separators.} We used vertex separators throughout. \emph{Edge separators} (delete edges to disconnect, à la sparsest cut, Chapter~33) are different and generally larger for the same balance; a bounded-degree graph converts between them up to a degree factor. Spectral methods (Cheeger, Chapter~32) find approximate \emph{edge} separators; planar/treewidth methods find exact \emph{vertex} separators.
  \item \textbf{Findability.} A separator theorem is only algorithmic if you can \emph{find} the separator fast. Planar separators: $\Oh(n)$. Treewidth-$w$ separators: need a decomposition, computable in $2^{\Oh(w)} n$ (Bodlaender) or approximately faster. General \emph{balanced} separator (sparsest cut) is NP-hard exactly, with an $\Oh(\sqrt{\log n})$-approximation (Arora--Rao--Vazirani, Chapter~33) --- the gap between ``exists'' and ``findable'' is itself a research frontier.
  \item \textbf{Recursiveness.} The class must be closed under the recursion. Planar minus a separator is still planar; treewidth-$w$ minus a bag is still $\le w$. General graphs offer no such closure, which is why the recursion does not start.
  \item \textbf{What the gain is.} Match the separator to the goal: $\Oh(\sqrt n)$ buys $\Oh(n\log n)$ D\&C, $2^{\Oh(\sqrt n)}$ exact, and $\Oh(n^{3/2})$ elimination; constant treewidth buys linear-time exact DP; sublinear-but-not-$\sqrt n$ separators (e.g.\ $\Oh(n^{2/3})$ for some bounded-genus or low-dimensional graphs) buy correspondingly weaker but still sub-exponential bounds.
\end{itemize}

\begin{trapbox}
A separator theorem with no efficient algorithm is a trap. Many graph classes \emph{have} small balanced separators that are NP-hard to find optimally (general graphs with small separators, via sparsest cut). Do not assume ``small separator exists'' implies ``divide-and-conquer is efficient'' --- you also need a fast (possibly approximate) separator oracle. Conversely, a fast oracle giving a \emph{slightly} larger separator (say $\Oh(\sqrt n\log n)$) usually costs you only a log factor in the recurrence and is well worth it. Optimize for \emph{found} size, not existential size.
\end{trapbox}

\begin{aestheticbox}
Step back and the whole chapter is one idea wearing three coats. \emph{Geometrically} (planar), a small separator is a short curve cutting a drawing in half. \emph{Structurally} (treewidth), it is a bag in a tree decomposition. \emph{Spectrally} (Chapters 32--34), it is a sparse cut revealed by the second Laplacian eigenvector. These are three windows onto ``where is this graph flimsy, and can I exploit the flimsiness recursively?'' Divide-and-conquer on graphs is the art of finding the flimsy spot, and a separator theorem is a guarantee that one exists and a recipe to find it. When it exists and you can find it, intractable problems become tractable; when it provably does not (expanders), the same problems stay hard --- and that dichotomy, separability versus expansion, organizes a huge swath of algorithmic graph theory.
\end{aestheticbox}

\section{Exercises}

\begin{exercise}[Grid lower bound]\label{ch27:ex-gridlb}
Prove that the $k\times k$ grid ($n=k^2$) has no $\tfrac12$-balanced separator of size smaller than $k=\sqrt n$. (Hint: a balanced separator must intersect $\Om(k)$ of the $k$ row-paths or $k$ column-paths; use Menger / vertex-disjoint paths to lower-bound the intersection.)
\end{exercise}

\begin{exercise}[Levels are separators]\label{ch27:ex-levels}
Implement BFS-level separators: given any graph, output, for each BFS level $L_i$, the sizes $\abs{L_{<i}},\abs{L_i},\abs{L_{>i}}$, and identify the median level. On a random planar graph, plot $\abs{L_i}$ versus $i$ and locate the thinnest level near the median. Where does balance fail, and where does size fail?
\end{exercise}

\begin{exercise}[Fundamental-cycle separator]\label{ch27:ex-cycle}
For a triangulated planar graph with a rooted spanning tree, write code that, for each non-tree edge, computes the inside-weight of its fundamental cycle, and returns the edge whose cycle is most balanced. Prove that the resulting cycle has both sides $\le\tfrac23 n$ and length $\Oh(h)$ where $h$ is the tree height.
\end{exercise}

\begin{exercise}[Break the median heuristic]\label{ch27:ex-break}
Construct a planar graph where the median BFS level has $\Th(n)$ vertices (so the ``just take the median level'' shortcut fails), forcing Movements 2--3 of the proof. (Hint: a graph whose BFS levels form a single fat level in the middle --- think of two large fans joined at a wide waist.)
\end{exercise}

\begin{exercise}[Treewidth $\Leftrightarrow$ separators]\label{ch27:ex-tw-sep}
Prove that if every subgraph of $\Gph$ has a $\tfrac23$-balanced separator of size $\le k$, then $\tw(\Gph)=\Oh(k\log n)$. (Build a tree decomposition top-down: separate, recurse, and add the separator to every descendant bag; bound the bag size by the recursion depth.)
\end{exercise}

\begin{exercise}[Sub-exponential vertex cover]\label{ch27:ex-vc}
Adapt Algorithm~\ref{ch27:alg-mis} to compute a minimum \emph{vertex cover} on planar graphs in $2^{\Oh(\sqrt n)}$ time, and verify the recurrence solves to sub-exponential. Then explain why the same approach gives a $3^{\Oh(\sqrt n)}$ algorithm for \emph{$3$-colouring}, with $3$ in place of $2$ per separator vertex.
\end{exercise}

\begin{exercise}[Nested dissection fill-in]\label{ch27:ex-nd}
On the $k\times k$ grid, order the vertices by nested dissection (separator last, recursively). Prove the fill-in obeys $F(n)\le 2F(n/2)+\Oh(n)=\Oh(n\log n)$ and the operation count is $\Oh(n^{3/2})$. Compare experimentally against a naive row-major ordering's fill-in.
\end{exercise}

\begin{exercise}[$r$-division by hand]\label{ch27:ex-rdiv}
On a $6\times 6$ grid, construct an $r$-division with $r=9$: partition into pieces of $\le 9$ vertices, each with $\Oh(\sqrt r)=\Oh(3)$ boundary vertices. Count the total boundary size and verify it is $\Oh(n/\sqrt r)$. Then answer a single shortest-path query between two corners using only the per-piece boundary distances.
\end{exercise}

\begin{exercise}[Expanders have no small separators]\label{ch27:ex-exp}
Using the (edge) expansion definition $h(\Gph)=\min_{\abs S\le n/2}\frac{e(S,\bar S)}{\abs S}\ge \beta$ for a constant $\beta>0$, prove that a constant-degree expander has every $\tfrac23$-balanced \emph{vertex} separator of size $\Om(n)$. Contrast with the planar theorem and explain why no $\Oh(\sqrt n)$-separator algorithm can exist for general bounded-degree graphs.
\end{exercise}

\begin{exercise}[Cut between two thin levels]\label{ch27:ex-twolevels}
In the planar proof we sometimes use two thin levels $L_a,L_b$ as the separator. Prove that $L_a\cup L_b$ is indeed a separator (no edge from $L_{<a}$ to $L_{>b}$ avoiding the band), and that if the middle band $M$ has $\abs M\le\tfrac23 n$ we are immediately done. Identify precisely which case forces the fundamental-cycle step instead.
\end{exercise}

\begin{exercise}[Trace the median-level argument]\label{ch27:ex-trace-bfs}
Reproduce Figure~\ref{ch27:fig-bfs} for the $6\times6$ grid ($n=36$) rooted at a corner. Write out the eleven level sizes $\abs{L_0},\dots,\abs{L_{10}}$, the cumulative counts, and identify the median level $i^*$. Show that $\abs{L_{i^*}}=6=\sqrt n$ and that both sides have $\le n/2$ vertices, so Movement~1 already separates the grid. Then change the root to a \emph{center} vertex and recompute: do the level sizes still peak at $\sqrt n$, and is the median level still thin?
\end{exercise}

\begin{exercise}[Treewidth DP table sizes]\label{ch27:ex-twdp-table}
Continue the worked DP trace (the labbox after Algorithm~\ref{ch27:alg-twdp}) on the graph of Figure~\ref{ch27:fig-td}, but now \emph{maximize a weighted} independent set with weights $w(a)=3$ and $w(b)=w(c)=w(d)=w(e)=1$. Recompute all three bag tables and the final answer. Then count: how many table entries does each bag store as a function of the width $w=2$, and how does that count grow if you instead solve \emph{dominating set}, where each bag vertex carries one of three states (in the set / dominated / undominated)? Relate the $3^{w+1}$ blow-up to the separator size.
\end{exercise}

\begin{exercise}[Edge vs.\ vertex separators on the grid]\label{ch27:ex-edge-vs-vertex}
The $k\times k$ grid has a vertex separator of size $k$ (a middle column) and also an \emph{edge} separator of size $k$ (the $k$ horizontal edges crossing the midline). Show both achieve $\tfrac12$-balance. Now take a bounded-degree graph of max degree $\Delta$ and prove that any edge separator of size $s$ yields a vertex separator of size $\le s$ (take one endpoint per cut edge), and that any vertex separator of size $t$ yields an edge separator of size $\le \Delta t$. Where does the grid sit in these bounds, and why is the conversion lossless there?
\end{exercise}
